% 21. Устойчивость против защит
\begin{frame}{Устойчивость против защит}
\framesubtitle{Защиты в пространстве признаков не спасают}
\vspace{1mm}
\begin{columns}[T,onlytextwidth]
\column{0.53\textwidth}
  \vspace{4mm}
  \includegraphics[width=\linewidth]{fig12.png}\par\vspace{3mm}
  \figcap{ASR против сертифицированной защиты BARS остаётся высоким по 8 датасетам.}
\column{0.44\textwidth}
\begin{wcard}[height=5.9cm]\scriptsize
  \cardtitle{Почему защиты не работают}\vspace{0.5mm}
  \iconitem[0.52cm]{badred}{ic_shieldv_w.png}{Adversarial training и randomized smoothing ограничивают норму в пространстве признаков}\vspace{3.2mm}
  \iconitem[0.52cm]{steel}{ic_arrows_w.png}{NetMasquerade правит трафик в физическом пространстве $\rightarrow$ это «неограниченная» атака}\vspace{3.2mm}
  \iconitem[0.52cm]{steel}{ic_bolt_w.png}{Малое возмущение трафика = большое расстояние в признаках, выходящее за сертиф. границу}\vspace{3.2mm}
  \iconitem[0.52cm]{accent}{ic_balance_w.png}{Иногда ASR даже выше: аугментация шумом снижает точность самой модели}
\end{wcard}
\end{columns}
\end{frame}
